% ============================================================================
% Section 5.2: Design Constraints  (FYDP-II)
%
% Paste into Chapter 5 after Section 5.1. It replaces the FYDP-I cost table
% (Table 5.2): Section 5.2.7 contains the updated table with actual costs.
% Citation keys: banglabert, mmbert, gemma4 (see new_references.tex).
% ============================================================================

\section{Design Constraints}
\label{sec:constraints}

This section describes the constraints that shaped the design and how the
implemented system meets them. Figures come from the measurements in
Chapter~\ref{ch:implementation} unless marked as estimates.

\subsection{Economic Constraint}

The project had no budget for cloud computing or paid APIs, so every component
had to run on hardware the team already owned and use free software. This was
met in full. All training, evaluation and inference ran on one consumer desktop
with an AMD RX 6600 GPU; the Google Colab subscription planned in FYDP-I was not
needed. All libraries are open source, and the pretrained models are free to
download.

The running cost of the finished system is close to zero. Checking an answer
takes 57~ms on an ordinary CPU and needs no GPU, no internet connection and no
per-request fee. The alternative the project set out to replace, an LLM used as
a judge, is far more expensive to run: on the same GPU, a local
12-billion-parameter judge needed 25.7 minutes to check the 588 test answers, where
BanglaBERT needed about 7 seconds, and a cloud LLM API would add a charge for
every request.

One economic limit concerns licensing. BanglaBERT is released under the Creative
Commons BY-NC-SA 4.0 licence~\cite{banglabert}, which allows research and
educational use but not commercial use. A commercial product would need
permission from its authors, or a backbone with a permissive licence. The mmBERT
comparison model is MIT-licensed~\cite{mmbert} and could be used commercially,
but it was slightly less accurate and too slow on a CPU (210~ms per answer).

\subsection{Environmental Constraint}

The system had to keep energy use low, both to train and, more importantly, in
everyday use, where it may check thousands of answers. A small model meets this
directly: one check uses 11.5~ms of GPU time, against 2,560~ms for the 12B LLM
judge, so the same work costs roughly 220 times less computation. On a CPU the
model needs no GPU at all.

Training was also modest. All experiments together, including the three-seed
runs, the comparison models and the LLM judges, took about 8 hours of machine
time on one desktop. At typical power draw for such a machine this is roughly
1 to 2~kWh (an estimate, not a measurement). LoRA reduced GPU training time per
epoch by a factor of 1.8. No new hardware was bought and no data-centre
resources were used, which avoids the embodied carbon of dedicated equipment.

\subsection{Ethical Constraint}

\paragraph{Privacy.} RAG systems often answer questions about health,
government services or personal matters. Because the detector runs entirely on
the local machine, the question, the retrieved context and the answer never
leave the device. No third-party service receives user data, unlike an
approach based on a cloud LLM API.

\paragraph{Transparency.} The system does not give an unexplained verdict. It
highlights the exact words it considers unsupported and shows a probability
for each word, so a user can check its reasoning against the context. The demo
also shows the lexical baseline beside the model.

\paragraph{Honest limits.} The detector is not perfect: it misses about 11\% of
hallucinated answers and wrongly flags about 4\% of faithful ones. It must be
presented as an aid to human judgement, not as a guarantee of correctness. The
results in this report are given with their baselines, with the variation
across three training runs, and with the system's weak types stated openly.

\paragraph{Data.} The dataset was built from publicly available sources: Bangla
Wikipedia (licensed CC BY-SA 4.0) and public Bangladeshi government portals. It
contains no personal data about private individuals. The pretrained models are
used within their licence terms, which for BanglaBERT means non-commercial
research use. Human annotation of the dataset is done by the team members
themselves.

\subsection{Health and Safety Constraint}

The system has no physical hazards; it is software running on standard office
computers. Its safety concern is informational. 500 rows of the dataset come
from the healthcare domain, including material from the Directorate General of
Health Services. A hallucinated medical answer that the detector fails to flag
could lead a user to act on wrong information, while a false alarm could make
a user doubt correct advice.

The design responds to this in three ways. The detector reports which words
are unsupported, so a reader can verify them rather than trust a single score.
The results are broken down by domain and hallucination type, so weak areas are
known before deployment. And the system is intended to support, not replace,
human review: in a health setting its output should be shown with a clear
notice that flagged and unflagged answers alike must be checked against an
authoritative source.

\subsection{Social Constraint}

Bangla is one of the most widely spoken languages in the world, yet it has few
natural language processing tools, and existing hallucination detectors and
datasets are almost all English. This project brings hallucination detection
to Bangla RAG systems, so that Bangla-speaking users of chatbots and question
answering services can be warned about unsupported answers.

The system was also designed to be affordable to the institutions most likely
to deploy it in Bangladesh, such as universities, small companies and
government offices, which often lack GPUs and budgets for cloud AI services.
It runs on an ordinary computer without a GPU. Highlighting the wrong words,
rather than giving a bare verdict, makes the output understandable to users
without technical training. A social risk remains: users may trust any answer
the detector does not flag. This is why the system presents itself as a
warning tool with known error rates.

\subsection{Political Constraint}

Two of the eight domains, Bangladesh Affairs and Government Services, include
political and administrative content such as elections, public office holders
and government procedures. Hallucinations on these subjects can spread
misinformation about public institutions. The detector is politically neutral
in one precise sense: it judges an answer only against the retrieved context
and has no opinion of its own on any political question. The consequence is
that it inherits the reliability of its source. If the retrieved context is
biased or wrong, an answer that repeats it will be judged supported. Deployers
must therefore choose trustworthy sources for retrieval.

Running locally also supports data sovereignty. Government bodies can check
answers without sending citizens' questions or internal documents to foreign
cloud providers, and the system keeps working without dependence on any
external service.

\subsection{Manufacturability and Cost Analysis}

The product is software, so manufacturability means that the system can be
reproduced, installed and run reliably on ordinary hardware. The code is
modular, with one configuration file choosing the model and training settings,
and its dependencies are listed for installation. Results are reproducible:
each configuration was trained with three fixed seeds and the variation between
them is reported. The data loader validates every dataset it reads and checks
that its word segmentation matches the labels, so a changed dataset fails
loudly instead of producing silently wrong results.

A deployment needs one of two artefacts: the fully fine-tuned model (420~MB) or
a LoRA adapter (3.4~MB) loaded on top of the public BanglaBERT base model. Any
computer with 16~GB of RAM can serve it through the FastAPI backend. At batch
size 1 this gives about 17 checks per second on the CPU used here and about 87
per second on the RX 6600 GPU.

Table~\ref{tab:cost} updates the FYDP-I cost estimate (Table~5.2) with the
actual costs of FYDP-II.

\begin{table}[htbp]
\centering
\small
\caption{Project cost analysis (actual, FYDP-II)}
\label{tab:cost}
\begin{tabular}{p{0.4cm}p{3.2cm}p{5.6cm}p{2.2cm}}
\hline
\textbf{No.} & \textbf{Item} & \textbf{Description} & \textbf{Cost (USD)} \\
\hline
1 & Development hardware & Existing laptops and one desktop with an AMD RX 6600 GPU & 0 (existing) \\
2 & Cloud GPU & Planned Google Colab Pro (about \$20) was not needed; all training ran locally & 0 \\
3 & Software & Python, PyTorch, Transformers, PEFT, FastAPI, Streamlit, llama.cpp & 0 (open source) \\
4 & Pretrained models & BanglaBERT (CC BY-NC-SA 4.0), mmBERT (MIT), Gemma~4 judges (Apache 2.0) & 0 \\
5 & Dataset & Built by the team from public sources & 0 \\
6 & LLM-judge comparison & Run locally instead of through a paid API & 0 \\
7 & Electricity & About 8 hours of machine time, roughly 1--2 kWh (estimate) & under 1 \\
8 & Documentation & Report printing, binding and presentation material & about 15 \\
\hline
 & \textbf{Total} & & \textbf{about 16} \\
\hline
\end{tabular}
\end{table}

The actual cost is below the FYDP-I estimate of about \$35, mainly because GPU
training on local hardware made the cloud subscription unnecessary. The cost of
using the system is also far below that of an LLM judge: it needs no API
subscription, and it checks an answer about 220 times faster than a local 12B
LLM on the same GPU.

\subsection{Sustainability}

The system is designed to stay useful and maintainable after the project ends.
\begin{itemize}
  \item \textbf{Cheap to retrain.} A full fine-tuning run takes 25 to 55 minutes
  on the RX 6600, and a LoRA run 13 to 29 minutes, so the model can be retrained
  whenever the dataset grows or is corrected by human annotation.
  \item \textbf{Cheap to extend.} Each LoRA adapter is 3.4~MB, so separate
  adapters for new domains can be added without storing another full model.
  \item \textbf{Open components.} All software is open source and all models
  have public weights, so the system does not depend on a vendor that could
  change its prices or withdraw its service.
  \item \textbf{Robust to data changes.} The dataset is still being extended.
  The pipeline measures sequence lengths and checks the data format on every
  load rather than assuming fixed dimensions, so new versions can be used
  without code changes.
\end{itemize}

Two risks to long-term maintenance remain. GPU training relies on DirectML,
which Microsoft now maintains only lightly and which requires an older PyTorch
release; a future move to a supported GPU platform may be needed. And
BanglaBERT reads at most 512 tokens, which may become a limit if future RAG
contexts are much longer.
